\subsection{5. Neuer kleiner Transportzeuge mit tatsächlichen
W-Kanälen}\label{neuer-kleiner-transportzeuge-mit-tatsuxe4chlichen-w-kanuxe4len}

Ein gemeinsamer Fermion-/Boson-Dreierstern hätte die einfache Identität

\[\adjustbox{max width=\linewidth}{$\displaystyle 
Q_x=b^\dagger s f_x,\quad Q_y=b^\dagger s f_y,
\qquad [Q_y^\dagger,Q_x]=(N_b+n_s)f_y^\dagger f_x.
$}\]

Auf festem \codeword{K=Nb+n\_s≥1} ist \codeword{c†=b†s/√K} exakt eine
CAR-Mode. Das ergibt eine gewöhnliche Dreimoden-Transferkette.
\textbf{Aber dieser Stern ist nicht in einer nativen W-Zeile enthalten:}
Jede Zeile ist ein Matching aus acht disjunkten Paaren. Die Analogie
allein wäre daher eine falsche Herkunftsangabe.

Im echten W gibt es stattdessen die Paare \codeword{(4,57)} in Kanal 0
und \codeword{(4,58)} in Kanal 1. Mit den sieben besetzten
Pauli-Blockern

\[\adjustbox{max width=\linewidth}{$\displaystyle 
S_b=\{8,16,28,32,44,52,56\}
$}\]

und einem \textbf{zusätzlichen} Mischer \codeword{J(b1†b0+b0†b1)}
schließt die volle native Wechselwirkung exakt auf den vier Zuständen

\[\adjustbox{max width=\linewidth}{$\displaystyle 
S_b+\{4,57\}\ \longleftrightarrow\ S_b+b_0
\ \longleftrightarrow\ S_b+b_1
\ \longleftrightarrow\ S_b+\{4,58\}.
$}\]

Alle haben N=9. Die vollständige Wirkung sämtlicher 480 signierter
Paarterme bleibt in diesem Raum, mit Matrix

\[\adjustbox{max width=\linewidth}{$\displaystyle 
H_4=\begin{pmatrix}
0&g&0&0\\g&\Delta&J&0\\0&J&\Delta&g\\0&0&g&0
\end{pmatrix}.
$}\]

Für \codeword{g/Δ=1/20}, \codeword{J/Δ=1-1/(10√3)} und
\codeword{tΔ=20π√3} gilt analytisch

\[\adjustbox{max width=\linewidth}{$\displaystyle 
P(57\to58)>\frac{2009992727}{2022609600}
=0.9937620819\ldots>99.3\%.
$}\]

Die Eigenfrequenzen Δ±J bleiben positiv. Dies ist keine unkontrollierte
Vierzustandsabschneidung: Zuerst ist die Invarianz unter der
vollständigen Wechselwirkung bewiesen, danach wird deren exakte
Einschränkung gelöst. Ohne Mischer zerfällt sie in zwei getrennte Blöcke
und der Transfer ist null.

Der reine Einzelmischer kommutiert nicht mit dem ursprünglichen Clock.
Die Clock-Orbitsumme der Mischer \codeword{(0,1)+(12,13)+(6,7)} ist
dagegen Clock-invariant und hat auf diesem Vierzustandsraum dieselbe
Wirkung. Somit wäre „jeder solche Transfer muss den Clock brechen``
falsch. Beide Mischer sind jedoch noch zusätzliche Operationen;
Clock-Invarianz allein beweist keine Erzeugbarkeit. Ein gemeinsamer
SU(4)-Cartan Q kommutiert mit H, Nb, Clock und beiden getrennten
Casimiren. Auf den vier Zuständen besitzt er aber die Werte
\codeword{(7,7,9,9)}. Die beiden Mischer ändern diese Quantenzahl:
\codeword{||[Q\_B,M01]||²\_F=8}, \codeword{||[Q\_B,Morb]||²\_F=24}.
\textbf{Auch die Clock-invariante Orbitsumme ist damit aus diesem
Operationssatz ausgeschlossen.}

Die noch fehlende Ressource ist hier konkret: eine Operation, die diese
innere Cartanladung verändert, oder eine global hergeleitete Kopplung an
einen ausdrücklich mitgerechneten Ladungsausgleich. Mehr Wörter aus dem
unveränderten erhaltenden Alphabet können beide Mischer nicht liefern.

Dieser Zeuge ist \textbf{keine} Propagation der
N=64-Grundzustands-Lochanregung, keine räumliche Zwei-Banken-Übertragung
und keine Konstruktion ihrer Präparation. Er zeigt konstruktiv, wie nah
eine kleine Paarumwandlung an einem Transfermechanismus liegen kann, und
benennt die fehlende Ressource.

\subsection{6. Das Feldwörterbuch: eine Basisänderung allein reicht
nicht}\label{das-feldwuxf6rterbuch-eine-basisuxe4nderung-allein-reicht-nicht}

Der bestätigte Grundfehler bleibt: \codeword{M\_A⊗ε\_Lorentz} ist
symmetrisch und verschwindet als Grassmann-Bilinear. Das betrifft die
lokale ableitungsfreie Zuordnung einer gleichhändigen Weylkopie je
ursprünglichem Label bei unverändertem antisymmetrischem W. Der
nichtverschwindende gleichhändige Kanal trägt \codeword{(1,0)}; daraus
folgt noch keine gesunde Kinetik.

Der neue Fünf-Zyklus im Trägergraphen widerlegt eine rein diagonale
Zuweisung entgegengesetzter Händigkeiten an jedes Paar. Ein neuer
eigener Satz geht darüber hinaus. Für irgendeine hermitesche interne
Händigkeit Γ müsste im unveränderten Ein-Kopie-Vektoransatz gelten

\[\adjustbox{max width=\linewidth}{$\displaystyle 
\Gamma^TM_A+M_A\Gamma=0\quad\forall A.
$}\]

Adjungieren und Einsetzen liefert \codeword{[Γ,M\_A†M\_B]=0}. Der
ursprüngliche Tensor faktorisiert exakt als

\[\adjustbox{max width=\linewidth}{$\displaystyle 
M_{ka}=S_k\otimes C_a,\quad
\sum_kS_k^\dagger S_k=5I_{16},\quad
\sum_aC_a^\dagger C_a=3I_4.
$}\]

Die S\_k†S\_l erzeugen M16; die C\_a†C\_b erzeugen M4. Dafür wurden
vollständige Rangzertifikate modulo 101 mit 256 bzw. 16 unabhängigen
ganzzahligen Matrixwörtern konstruiert. Ein nichtverschwindender
Determinant modulo 101 beweist Nichtverschwindung über C; keine
Rangdefizienz wird übertragen. Damit ist Γ skalar, und die ursprüngliche
Gleichung erzwingt Γ=0.

\textbf{Keine hermitesche Involution Γ²=I auf demselben Träger repariert
diesen Ein-Kopie-Vektoransatz.} Der Ausschluss ist nun basisunabhängig,
aber weiterhin vertragsgebunden.

Nicht ausgeschlossen sind zusätzliche Dirackomponenten, ein unabhängiger
Zweierfaktor oder Ableitungsterme. Insbesondere ist

\[\adjustbox{max width=\linewidth}{$\displaystyle 
M_A\otimes\epsilon_{aux}\otimes\epsilon_{Lorentz}
$}\]

antisymmetrisch und nicht null. Ein ableitungshaltiger gleichgeladener
Vektorstrom \codeword{M\_IJ ε\_ab ψ\_Ia ↔∂\_μ ψ\_Jb} ist ebenfalls nicht
null; der pauschale Vektorausschluss durch Ladung gilt nicht außerhalb
der ableitungsfreien Klasse.

Die Zahlen 180/128 oder 60/256 zählen Komponenten eines bestimmten
Feldansatzes, nicht automatisch unabhängige propagierende Oszillatoren.
Schur fixiert interne Intertwiner, nicht alle Lorentz-/Ableitungsterme.
Der behauptete Zwang „zuerst räumlich skalieren, danach Feldtyp prüfen``
folgt nicht daraus. Zuerst bleibt ein kleiner konsistenter
Typ-/Kinetiktest sinnvoll; seine reale räumliche Herkunft muss
anschließend gemeinsam geprüft werden.

\subsection{7. Was der nachgereichte Universalraum-Text
trägt}\label{was-der-nachgereichte-universalraum-text-truxe4gt}

\subsubsection{7.1 Eine sinnvolle Hypothese, kein bereits
identifiziertes universelles
Objekt}\label{eine-sinnvolle-hypothese-kein-bereits-identifiziertes-universelles-objekt}

Nichtorthogonale oder überlappende lokale Einbettungen in eine
gemeinsame Algebra sind eine ernstzunehmende Alternative zu unabhängigen
Banken. Sie können Voraussetzungen der lokalen Paritätsschranke ändern.
Das ist ein Forschungsauftrag, kein Widerspruch zur bisherigen Schranke.

Bei Einbettungen \codeword{J\_A:V\_A→V} gilt für die CAR jedoch

\[\adjustbox{max width=\linewidth}{$\displaystyle 
\{f_A(u),f_B(v)^\dagger\}=\langle J_Au,J_Bv\rangle.
$}\]

Die Überlappungs-Grammatrix gehört also zwingend zum Modell. Die alten
Produktzustände \codeword{Ω\_A⊗Ω\_B}, ihre beiden unabhängigen
Ladungssektoren und der Zwei-Banken-Polraum dürfen nicht ungeprüft in
einen überlappenden Träger übernommen werden. Der gemeinsame
Hamiltonoperator und sein Zustand müssen dort neu beziehungsweise durch
einen bewiesenen Adapter konstruiert werden.

\subsubsection{7.2 Der vorgeschlagene Kill-Test muss korrigiert
werden}\label{der-vorgeschlagene-kill-test-muss-korrigiert-werden}

Für unabhängige orthogonale Banken kann Π\_AΠ\_B deren gemeinsame
Parität sein. Für überlappende Unterräume gilt das nicht allgemein; ihre
lokalen Paritäten müssen nicht einmal kommutieren. Der Test
\codeword{[U,Π\_AΠ\_B]=0} ist deshalb im neuen Bild nicht der richtige
universelle Test. Stattdessen muss die tatsächliche globale Parität
Π\_global separat definiert und erhalten werden. Ein
nichtverschwindender Kommutator mit einer lokalen Parität belegt zudem
allein noch keinen gerichteten Transfer.

\subsubsection{7.3 Der einfachste False-Positive: Ein Signal ohne
Bewegung}\label{der-einfachste-false-positive-ein-signal-ohne-bewegung}

Seien \codeword{a=c1}, \codeword{b=(c1+c2)/√2} auf zwei globalen
CAR-Moden und \codeword{H=ωN}. Dann haben die beiden
Einteilchen-Chartzustände bereits Überlappung \codeword{1/√2}. In ihrer
nichtorthogonalen Basis ist \codeword{K=ωS} mit einem
nichtverschwindenden Offdiagonalelement. Trotzdem bleibt die betreffende
Wahrscheinlichkeit zu jeder Zeit genau 1/2. Es wurde nichts
transportiert.

Ein belastbarer gemeinsamer Test muss daher mindestens

\[\adjustbox{max width=\linewidth}{$\displaystyle 
S_{AB}=\langle\psi_A,\psi_B\rangle,\qquad
K_{AB}=\langle\psi_A,H\psi_B\rangle
$}\]

und die tatsächliche Zeitantwort gemeinsam bestimmen. Auf einem positiv
definiten Gramträger ist
\codeword{S\textasciicircum{}\{-1/2\}KS\textasciicircum{}\{-1/2\}} der
korrekte orthonormierte Operator; bei singulärem S ist zunächst der
Nullraum zu quotientieren. Im Gegenbeispiel ergibt dies nur ωI. Die
Offdiagonale von K war kein Hop.

Das trifft sogar unmittelbar den vorhandenen nativen Pol. Dessen
64-dimensionaler Raum besitzt genau \textbf{eine} Energie E\_h, also
\codeword{P\_h H P\_h=E\_h P\_h}. Sind A und B nur neue Charts innerhalb
dieses selben Polraums, gilt zwangsläufig \codeword{K=E\_h S}. Bloßes
Umbenennen der 64 entarteten Polrichtungen erzeugt daher keine
Ausbreitung durch H. Eine tatsächliche globale Dynamik beziehungsweise
eine verfügbare aktive Operation muss hinzukommen und aus der Quelle
ausgewiesen werden.

Für die eigentliche TFPT-Lochantwort wäre der gemeinsame Gegenstand

\[\adjustbox{max width=\linewidth}{$\displaystyle 
G_{AB}(t)=\langle\Omega,
 f_B^\dagger e^{-it(H-E_0)}f_A\Omega\rangle,
$}\]

mit einer einzigen Quelle, einem H und einem Ω. \codeword{G(0)} bezahlt
die schon vorhandene Überlappung; die dynamische Änderung muss davon
unterschieden werden. Das ist die geschärfte, kleinere nächste
Forschungsaufgabe.

\subsubsection{7.4 Was nicht automatisch
folgt}\label{was-nicht-automatisch-folgt}

\begin{itemize}
\tightlist
\item
  \textbf{Zeit:} Für einen Grundzustand ist
  \codeword{e\textasciicircum{}\{-itH\}Ω=e\textasciicircum{}\{-itE0\}Ω};
  sein physischer Zustand bleibt gleich. Das ist kein eigener Zeitpfeil.
  Nichtstationäre Präparation, Korrelationen und Records müssen
  ausgewiesen werden.
\item
  \textbf{Eichfeld:} Für bloße Basiswechsel \codeword{Uxy=Vx†Vy}
  teleskopiert jede geschlossene Holonomie zu I. Echte Krümmung braucht
  zusätzliche Verbindungs-/Transportdaten oder einen nachgewiesenen
  Mechanismus.
\item
  \textbf{Spin:} Zwei Orientierungsmarken liefern nicht von selbst einen
  SL(2,C)-Spinor, dessen Transformationsgesetz und Kinetik. Der
  zusätzliche antisymmetrische Zweierfaktor muss unabhängig konstruiert
  sein.
\item
  \textbf{Raum:} Minimale Operationskosten sind ohne Reversibilität
  zunächst gerichtete Kosten, nicht notwendig eine symmetrische Metrik.
  Kubisches Volumenwachstum allein beweist keine 3D-Mannigfaltigkeit,
  keine Lorentzstruktur und keinen gemeinsamen Lichtkegel.
\item
  \textbf{Gravitation:} Dynamische Beziehungen sind ein möglicher
  Träger, aber ein masseloser Spin-2-Sektor, zwei Helizitäten,
  Energiepositivität und universelle konsistente Kopplung bleiben zu
  beweisen.
\item
  \textbf{Arithmetik:} Primitive Schleifen besitzen generisch neue
  gemischte primitive Zyklen. Sie sind nicht automatisch Primzahlen. Die
  Xi-Determinante bleibt eine unbewiesene vollständige
  Operatoridentität. Der getrennte RH-Anhang dokumentiert Vorarbeiten,
  Beispiele und Quellenstatus.
\end{itemize}

Der Text wird deshalb als \textbf{präzisierter globaler
Forschungsansatz} integriert, nicht als Nachweis, dass derselbe
unbekannte Baustein schon alle T1-T8-, RH-, Faktorisierungs-, P/NP- und
Hylæan-Fragen beantwortet.

\subsection{8. T1-T8: gemeinsamer Fortschritt ohne
Statusinflation}\label{t1-t8-gemeinsamer-fortschritt-ohne-statusinflation}

{\def\LTcaptype{none} % do not increment counter
\begin{longtable}[]{@{}
  >{\raggedright\arraybackslash}p{(\linewidth - 4\tabcolsep) * \real{0.3333}}
  >{\raggedright\arraybackslash}p{(\linewidth - 4\tabcolsep) * \real{0.3333}}
  >{\raggedright\arraybackslash}p{(\linewidth - 4\tabcolsep) * \real{0.3333}}@{}}
\toprule\noalign{}
\begin{minipage}[b]{\linewidth}\raggedright
Front
\end{minipage} & \begin{minipage}[b]{\linewidth}\raggedright
In dieser Revision weitergekommen
\end{minipage} & \begin{minipage}[b]{\linewidth}\raggedright
Noch erforderlicher Nachweis
\end{minipage} \\
\midrule\noalign{}
\endhead
\bottomrule\noalign{}
\endlastfoot
T1 - Ursprung und Auswahl & Explizite Identität zwischen Quell-Clock und
innerer Spinwirkung; Operationsverträge genauer & Primitive Auswahl von
P1/P2, Träger, erlaubten Instrumenten und Zustand \\
T2 - Half-Charge/E8-Feld & Bestehender geladener Pol bleibt abgesichert,
innere Darstellung präzisiert & Renormiertes Half-Charge-Feld mit
Energie-/Adjungiertenkontrolle und E8-/Clock-Adapter \\
T3 - gemeinsamer 3+1D-Parent & Kleiner exakter Paartransferzeuge und
korrigierter Überlappungstest & Quellenseitig ausgewählter räumlicher
Parent auf einem gemeinsamen Hilbertraum und Zustand \\
T4 - chirale Materie & Basisunabhängiger Ausschluss eines konkreten
Ein-Kopie-Vektoradapters & Konsistenter chiraler Feldadapter,
Maß/Anomalien/Index und Spiegelentkopplung \\
T5 - Kontinuum und Dynamik & Exakte endliche Clock-/Casimir-Kombination
und bedingter invarianten Transferblock & Kontrollierter
wechselwirkender Grenzübergang, Lorentzverhalten, Clustering und
Streuung \\
T6 - Kopplungen und Texturen & Keine neue Herleitung; neue
Mischparameter sichtbar zusätzlich & Herkunft aller Kopplungen und
vollständiger Neutrinostruktur/-skala \\
T7 - Gravitation & Der falsche pauschale Tensor-Ausschluss bleibt
zurückgenommen & Masseloser dynamischer Spin zwei, beide Helizitäten,
konsistente universelle Kopplung im selben Parent \\
T8 - Präparation und Auslesung & Gram-/Antworttest gemeinsam formuliert;
spezielle Blockerressource konkret & Natives geladenes Instrument,
Quellzustandswahl, Records und ein gemeinsames Auslesefunktional \\
\end{longtable}
}

Keiner der acht vollständigen Abschlussverträge wird in dieser Revision
auf „gelöst`` gesetzt. Gelöst wurden ausdrücklich benannte Teilfragen.

\subsection{9. Die nächsten drei Arbeiten - klein, gemeinsam und
entscheidbar}\label{die-nuxe4chsten-drei-arbeiten---klein-gemeinsam-und-entscheidbar}

\textbf{A. Den gemeinsamen Quellenvertrag wirklich hinschreiben.} Zwei
lokale Einbettungen, ihr CAR-Gram, globale Parität und Ladung,
ursprüngliche Operationen und ein einziges H angeben. Abbruch für eine
Kandidatenfassung, wenn sie nur zwei Koordinatenansichten desselben
unveränderten Signals umbenennt. Die gemeinsamen Generatoren müssen
einen nachweislich neuen Antwortprozess ermöglichen, nicht nur eine
andere Beschriftung.

\textbf{B. Die fehlende Zwischenoperation ausweisen.} Der
Vierzustandszeuge gibt eine konkrete positive Zieloperation vor. Zu
prüfen ist, ob der Compiler ein geeignetes Mischinstrument oder eine
andere geteilte Paarstruktur wirklich liefert. Clock und Casimire
alleine sind kein Herkunftsbeweis. Zeigt ein erhaltener gemeinsamer
Generator die Unmöglichkeit für einen festen Operationssatz, endet
dessen weitere Kommutator-Suche; dann muss ein anderer bereits
vorhandener Quellbestandteil benannt werden.

\textbf{C. Dieselbe Anregung und denselben Zustand behalten.} Auf dem
Kandidaten \codeword{G\_AB(0)} und \codeword{G\_AB(t)} einschließlich
vollständiger Restantwort bestimmen. Der N=9-Zeuge darf nicht still zum
N=64-Lochpol umgedeutet werden. Ein minimaler Lorentz-/Kinetikadapter
muss die nachgerechneten Tensorbedingungen erfüllen. Erst mit diesem
gemeinsamen Anschluss trägt eine großräumige Skalierungsrechnung.

Der stärkste Vereinfachungsgewinn dieser Runde ist damit konkret:
\textbf{Clock und innere Symmetrie sind nicht zwei voneinander
unabhängige Mechanismen; reine Überlappung und echte Bewegung sind
dagegen zwei verschiedene Dinge.} Die weitere Suche sollte diese erste
Identität nutzen und die zweite Unterscheidung im selben kleinen Modell
erzwingen.

\subsection{10. Belege und
Reproduzierbarkeit}\label{belege-und-reproduzierbarkeit}

Die drei Eingangstexte bleiben unverändert im Prüfpaket. Ein Teil-Audit
reproduziert die 1127 gelieferten Guards; eigene Programme behandeln
Clock/Casimir, Quellenreichweite, direkte Krylov-Kontraktion,
basisunabhängige Feldbedingungen, den tatsächlichen Vierzustandskanal
und die Überlappungsgegenkontrollen. Exakte und numerische Teile sind in
den jeweiligen Ergebnisdateien getrennt markiert. Die abschließenden
Zählungen und Bytevergleiche stehen im Auslieferungs-/Replayprotokoll;
interne erneute Aufrufe der 944 Guards werden nicht mehrfach als neue
Tests gezählt.

Die drei ausführlichen parallelen Berichte und der
Überlappungs-/RH-Nachtrag folgen im vollständigen Dokument. Ihr Status
ersetzt widersprechende historische Deutungen, nicht unveränderte ältere
Beweise. Das native frühere Archiv bleibt zusätzlich unverändert im
Paket; dessen große frühere Enumeration wird nicht als in dieser
Revision erneut durchgeführt ausgegeben.

\begin{center}\rule{0.5\linewidth}{0.5pt}\end{center}

